\section{統計的検定と推定 : 数学的背景と応用}
  \subsection{一致性と不偏性}

    Part Aでは, 標本から検定統計量を計算し, 帰無仮説のもとでの珍しさを測って判断を下す, という検定の流れを追った.
    Part Bでは, その中身を数式と数値で確かめていく.

    検定で本当に比べたいのは, 手元の標本そのものではなく, その背後にある母集団どうしである.
    たとえばサプリの効果を知りたいなら, サプリを飲んだ人の全体と飲んでいない人の全体を比べられれば, 話はそこで終わる.
    しかし, そんな大きな集団のデータを丸ごと手に入れるのは現実的ではない.
    そこで, 母集団の情報（母数）を標本からもっともらしく推定し, その推定値どうしを比べることになる.

    すると今度は, 推定のもっともらしさをどう評価するかが問題になる.
    その観点として, 5章の終わりのコラムでも触れた二つの性質がある.

    \begin{itemize}
      \item \textbf{一致性：} サンプルサイズを大きくしていくと, 推定量が母数に限りなく近づいていく.
      \item \textbf{不偏性：} 推定量の期待値が母数と一致する.
    \end{itemize}

    ぱっと見にはよく似ていて, 同じことを言っているようにも見える.
    しかし, 別の性質である\footnote{たとえば, サンプルサイズが大きくなってもノイズの影響が消えないような推定量では, 不偏性はあるのに一致性はない, ということも起こり得る.}.
    一致性はとても自然な要求だが, 一致性だけでは, 莫大な数のデータを集めないと母数に十分近づかない場合もある.
    そこを補うのが不偏性である.
    推定量の期待値が母数と一致するなら, 推定の狙いは平均的には高くも低くも偏っていない.
    サンプルサイズが限られた手元のデータからの推定でも, 系統的なずれを心配せずに扱える.
    このように一致性と不偏性は互いに補い合う性質で, 検定に使う推定量には両方を備えたものを選ぶのが基本になる.
    標本平均はその代表例で, 一致性と不偏性の両方を持っている.


  \subsection{自由度}

    本章の検定には, この後\textbf{自由度}という量が繰り返し登場する.
    t分布の形は自由度で決まり, カイ二乗分布の形も自由度で決まる.
    つまり, 自由度を取り違えると, 同じ検定統計量からでも違う$p$値が出てしまう.
    $p$値に直結する量でありながら, 自由度は「データの数から1を引いたもの」という公式の形で丸暗記されがちで, 意味を問われると答えに困る量でもある.

    幸い, 自由度の意味が最もはっきり見える場面が一つある.
    母集団の分散を標本から推定する場面だ.
    そこから話を始めよう.

    \subsubsection*{母集団の分散を推定したい. ところが母平均がわからない}

    母集団の分散（母分散） $\sigma^2$ は, 母集団の各データが母平均 $\mu$ からどれだけずれているかの二乗を, 母集団全体で平均した量である.
    各データのずれを二乗して足し, データ数で割る.
    式で書けば $\sigma^2 = \dfrac{1}{N}\sum (X_i - \mu)^2$ のようになるが, 構造としては「平均からのずれの二乗の平均」と覚えておけば十分だ.

    では, 手元の標本からこの量を推定したい.
    ここで困った事態に直面する.
    定義の中身に登場する $\mu$ がわからない.
    そもそも母集団のことが直接わからないからこそ, 標本を使って推定しようとしているのだった.
    それなのに, 推定したい量の定義の中に, もう一つの「わからない」量が入っている.

    打開策はひとつしかない.
    手元の標本から計算できる量で $\mu$ を代用する.
    最も自然な代用品は, 標本平均 $\bar{X} = \frac{1}{n}\sum X_i$ だ.
    そこで $\mu$ を $\bar{X}$ で置き換えた量
    \[
      \frac{1}{n} \sum_{i=1}^n (X_i - \bar{X})^2
    \]
    を考えてみたくなる.
    第2章で定義した分散と同じ形だし, これでうまくいきそうな気もする.
    しかし, この素直な置き換えには, ひとつ落とし穴がある.

    \subsubsection*{\texorpdfstring{$\bar{X}$}{X̄} で置き換えると, 何がずれるのか}

    $\bar{X}$ で代用したとき, 本当に知りたい量とどれくらい差が出るのかを調べてみよう.
    母分散の分子に当たる $\sum(X_i - \mu)^2$ と, 標本から計算する $\sum(X_i - \bar{X})^2$ の差を計算する.

    手始めに, $X_i - \mu = (X_i - \bar{X}) + (\bar{X} - \mu)$ と分解してから二乗を展開する：
    \begin{align*}
      \sum_{i=1}^n (X_i - \mu)^2
      &= \sum_{i=1}^n \left[ (X_i - \bar{X}) + (\bar{X} - \mu) \right]^2 \\
      &= \sum_{i=1}^n (X_i - \bar{X})^2 + 2(\bar{X} - \mu) \sum_{i=1}^n (X_i - \bar{X}) + \sum_{i=1}^n (\bar{X} - \mu)^2
    \end{align*}
    第二項の和の部分 $\sum (X_i - \bar{X})$ は, $\bar{X}$ の定義から常にゼロになる.
    よって第二項そのものが消える.
    結果として
    \[
      \sum_{i=1}^n (X_i - \mu)^2 = \sum_{i=1}^n (X_i - \bar{X})^2 + n(\bar{X} - \mu)^2
    \]
    が得られる.

    左辺は, 各データが\textbf{母平均から}どれだけずれているかの総量で, 本当に知りたい量である.
    第一項は, 各データが\textbf{標本平均から}どれだけずれているかの総量で, こちらは手元で計算できる.
    残る第二項は, 代用品の $\bar{X}$ 自身が本物の $\mu$ からどれだけずれているかを, データ全体分（$n$倍）まとめた量だ.

    つまり, 本物の $\mu$ から測ったずれの総量と, 代用品の $\bar{X}$ から測ったずれの総量は, 「$\bar{X}$ 自身のずれ」の分だけ食い違う.
    分解の式で, その食い違いが第二項として取り出せた.
    では, この食い違いは, どのくらいの大きさなのか.
    期待値を取って測ってみよう.

    \subsubsection*{ずれの大きさはデータ何個分にあたるか}

    上の式の各項の期待値を求めて, 大きさを比べる.

    左辺の期待値は, 母分散の定義から各 $X_i$ について $E[(X_i - \mu)^2] = \sigma^2$ なので
    \[
      E\!\left[\sum_{i=1}^n (X_i - \mu)^2\right] = n\sigma^2
    \]
    である.
    各データが「母平均から $\sigma^2$ 分ずれる」のを $n$個足したぶん, ということだ.

    第二項はどうか.
    4章で見たように, 標本平均 $\bar{X}$ 自体が母平均からどれだけずれるかは標準誤差で測れて, その分散は $\sigma^2/n$ だった.
    つまり $E[(\bar{X} - \mu)^2] = \sigma^2/n$.
    したがって
    \[
      E[n(\bar{X} - \mu)^2] = n \cdot \frac{\sigma^2}{n} = \sigma^2
    \]
    となる.

    第二項の期待値は $\sigma^2$ ちょうど, データひとつ分の分散ぴったりである.
    偶然ではなく, 構造的な必然だ.

    どういうことか.
    標本平均 $\bar{X}$ は, 母平均 $\mu$ を推定するために, $n$個のデータの中心の情報を一つの数に集約した量だ.
    そこまでは順当である.
    問題は, この $\bar{X}$ を, さらに分散の推定にも使い回そうとした点にある.
    本当に欲しいのは, 各データが母平均 $\mu$ からどれだけずれているかの総量（左辺）だった.
    しかし手元で計算できるのは, 各データが標本平均 $\bar{X}$ からどれだけずれているかの総量（第一項）である.
    平均の推定と分散の推定を, どちらも同じ $n$個のデータで済まそうとしたために, 中心を $\bar{X}$ で代用した代償が生じる.
    その代償の大きさが, 期待値で測るとちょうどデータひとつ分の分散 $\sigma^2$ なのだ.

    \subsubsection*{自由度 \texorpdfstring{$n-1$}{n-1} の意味}

    左辺の期待値が $n\sigma^2$, 第二項の期待値が $\sigma^2$ なので, 第一項の期待値は引き算で
    \[
      E\!\left[\sum_{i=1}^n (X_i - \bar{X})^2\right] = n\sigma^2 - \sigma^2 = (n-1)\sigma^2
    \]
    と求まる.

    手元で計算できる $\sum(X_i - \bar{X})^2$ が持っているずれの情報は, 期待値でみると, 本当に欲しかった $n\sigma^2$ よりデータひとつ分だけ少ない.
    だから, この量を $n$ で割ると, 母分散より少し小さめの値が出やすい.
    一方, $n-1$ で割れば, 期待値はちょうど $\sigma^2$ に一致する.
    平均的には母分散をぴったり言い当てる量になるわけだ.

    この $n-1$ で割った量を\textbf{不偏分散}と呼ぶ：
    \[
      s^2 = \frac{1}{n-1} \sum_{i=1}^n (X_i - \bar{X})^2
    \]
    \textbf{不偏}とは, 期待値を取ったときに偏りなく母数に一致する性質のことで, 前節の不偏性がこれにあたる.
    また, 大数の法則から, $n$ を大きくすれば標本平均 $\bar{X}$ は母平均 $\mu$ に近づくので, 不偏分散も $n$ を大きくすれば母分散に近づく.
    つまり一致性も併せ持つ
    \footnote{
      第2章で定義した分散は $n$ で割るものだった.
      あちらは手元のデータの散らばりをそのまま要約する量, こちらは母分散を狙う推定量で, 役割が違う.
      $n$ で割る版は不偏ではないが, 一致性は持つ.
      第2章の脚注で予告した「$n-1$ で割る理由」の答えが, この節の計算である.
    }.

    手元の $n$個のデータから出発して, $\bar{X}$ という一つの代表値に集約する過程で, 使える情報がデータひとつ分だけ目減りした.
    残った情報の個数が, データの数 $n$ から一つ引いた $n-1$ である.
    この数を\textbf{自由度}と呼ぶ.
    分散を $n-1$ で割って不偏にする補正は, 目減り分の埋め合わせだ.
    集約の代償として失われた一個分を, 分母を一つ減らすことで取り戻している.

    極端な場合で確かめておこう.
    データが1個しかないとき, 標本平均はそのデータ自身になり, ずれの総量 $\sum(X_i - \bar{X})^2$ は必ずゼロになる.
    ゼロを $n = 1$ で割れば「分散はゼロ」, つまり「母集団はまったくばらついていない」という見積もりが出てくるが, データ1個でばらつきの話ができるはずがない.
    一方, $n - 1 = 0$ で割ろうとすると, 値は定まらない.
    ばらつきの情報が何も残っていないから見積もり不能, という答えであり, こちらのほうが実態に合っている.
    平均の推定に1個分の情報を使い切ると, ばらつきに回せる情報はゼロになる.
    自由度ゼロとは, そういう状態のことだ.

    \subsubsection*{もう一つの見方: 標本平均を固定して残りの自由さを見る}

    自由度 $n-1$ を, もう一つ別の角度から眺めてみたい.
    この後のt検定やカイ二乗検定に出てくる自由度も, 同じ角度から数えられる.

    標本から $\bar{X}$ を計算した.
    ここで, $\bar{X}$ の値を一旦固定する.
    そのうえで, 「同じ標本平均をもつ別のデータの組」を考えてみる.
    つまり, $\frac{1}{n}\sum X_i = \bar{X}$ という条件を満たしたまま, $X_1, \ldots, X_n$ を動かしてみるということだ.

    この条件のもとでは, $X_1, \ldots, X_{n-1}$ までは自由に動かせる.
    しかし最後の $X_n$ は
    \[
      X_n = n\bar{X} - \sum_{i=1}^{n-1} X_i
    \]
    と自動的に決まってしまう.
    平均を $\bar{X}$ に合わせるためには, 最後の一つで帳尻を合わせるしかないからだ.
    自由に動かせる個数は $n-1$ にとどまる.

    この $n-1$ という数は, 元のデータそのものに最初から制約があるという話ではない.
    元のデータは $n$個とも自由に観測される.
    しかし, そこから $\bar{X}$ という代表値を計算して固定したうえで, 残ったずれ $(X_i - \bar{X})$ を眺めると, ずれの組が自由に動ける余地は $n-1$個分しかない.
    平均を引いたあとのずれに残された, 自由さの個数.
    「自由度」という名前は, この見方から来ている.

    この見方と, 前段の「集約の代償」は, 同じ事情の別の表れである.
    $\bar{X}$ という代表値を作って固定したことが, 期待値の計算では「データひとつ分の分散が代償として現れる」形で見え, ずれの動きに注目すると「自由に動ける個数が一つ少ない」形で見える.

    \subsubsection*{自由度は分散の推定だけの話ではない}

    ここまでは分散の推定を舞台にしてきたが, 同じ構造は統計的検定のあちこちにある.
    標本から推定量を一つ計算するたびに, 情報がその推定量に集約され, 自由に残る情報が一つ減る.

    \begin{itemize}
      \item \textbf{t分布}: 標本平均と標本分散を組み合わせて作るため, 分散の自由度 $n-1$ が分布の形を決める.
      \item \textbf{カイ二乗分布}: 正規変数の二乗和として作られるが, 推定に使った個数分だけ自由度が減る.
      \item \textbf{F分布}: 二つの不偏分散の比を取るため, 分子と分母それぞれの自由度が分布の形を決めるパラメータになる.
    \end{itemize}

    どの場合も構造は共通で, 推定に使った個数分だけ自由に残る情報が目減りし, その目減りを反映した数が自由度になる.
    そして次の節から見るように, 自由度はt分布やカイ二乗分布の形を決め, $p$値の値を左右する.
    自由度を取り違えることは, 珍しさを比べる先の分布を取り違えることにあたる.


  \subsection{t検定の数理構造と計算の実際}

    Part Aでは, 原作あり群と原作なし群の平均視聴回数に差があるかという問いを, t検定の手順に載せるところまで進めた.
    ここでは同じ例に実際の数値を入れて, 検定統計量の計算から$p$値の読みまでを一段ずつ確かめる.

    使うのは, 5章の実践で差に幅をつけたのと同じ2群のデータである.
    \begin{itemize}
      \item \textbf{原作あり群}（37本）：平均視聴回数 $= 274.1$万回, 標準偏差 $= 257.1$万回
      \item \textbf{原作なし群}（43本）：平均視聴回数 $= 153.7$万回, 標準偏差 $= 207.2$万回
    \end{itemize}
    このデータから, 「原作あり群のほうが平均視聴回数が高い」と言えるかどうかを検定したい.

    \subsubsection*{仮説を立てる}

    検定の出発点は仮説の設定だった.
    \begin{itemize}
      \item 帰無仮説 $H_0$：原作の有無で平均視聴回数に差はない（$\mu_1 = \mu_2$）
      \item 対立仮説 $H_1$：原作あり群の平均視聴回数のほうが高い（$\mu_1 > \mu_2$）
    \end{itemize}
    特定の方向に注目しているので, 片側検定を行うことになる.

    \subsubsection*{t統計量の定義と計算}

    t検定では, 二つの標本平均の差が, ばらつきに比べてどれだけ大きいかを測る.
    各群のばらつきが違いうることを認める立場で, 次の統計量を用いる：

    \[
      t = \frac{\bar{x}_1 - \bar{x}_2}{\sqrt{\dfrac{s_1^2}{n_1} + \dfrac{s_2^2}{n_2}}}
    \]

    分母は, 5章のケース3で差に幅をつけたときに登場した式と同じものだ.
    各群の分散をそれぞれの件数で割って足し合わせ, 平方根をとっている.
    分子の「差そのもの」を, 分母の「差のばらつきの大きさ」で割ることで, ばらつきを基準に差の大きさを測り直している.

    この式は, $s_1$ と $s_2$ を最後まで別々に扱っている.
    両群の母分散が等しいと仮定すれば, 二つの標本分散をひとまとめにして扱う書き方もできる.
    その立場で組まれたのが古典的なt検定（Studentのt検定）だ.
    一方, 本書で採用するのは, 両群の分散が違いうるという立場をとる\textbf{Welchのt検定}で, その分母として上の式が現れている.
    両者の関係は, 次の節で改めて整理する.

    数値を代入して計算してみよう.

    \begin{itemize}
      \item $n_1 = 37,\quad n_2 = 43$
      \item $s_1 = 257.1,\quad s_2 = 207.2$
      \item $\bar{x}_1 = 274.1,\quad \bar{x}_2 = 153.7$
    \end{itemize}

    まず, 分母の中身を求める：

    \[
      \frac{s_1^2}{n_1} + \frac{s_2^2}{n_2} = \frac{66100.41}{37} + \frac{42931.84}{43} = 1786.50 + 998.41 = 2784.91
    \]

    その平方根が分母になる：

    \[
      \sqrt{2784.91} \approx 52.77
    \]

    したがってt統計量は

    \[
      t = \frac{274.1 - 153.7}{52.77} \approx \frac{120.4}{52.77} \approx 2.28
    \]

    \subsubsection*{自由度と\texorpdfstring{$p$}{p}値}

    t分布を使って$p$値を求めるには, 自由度を決める必要がある.
    Welchのt検定では, 自由度を次の式で計算する：

    \[
      \nu = \frac{\left( \dfrac{s_1^2}{n_1} + \dfrac{s_2^2}{n_2} \right)^2}
                  {\dfrac{(s_1^2/n_1)^2}{n_1 - 1} + \dfrac{(s_2^2/n_2)^2}{n_2 - 1}}
    \]

    一見いかつい式だ.
    なぜこんな式になるのか, そして「自由度」と呼んでいるのに整数ですらない値が出てくるのはどういうことなのか.
    次の節でじっくり扱うことにして, ここでは, この式に従って計算するとどうなるかだけ確認しておく.

    分子：
    \[
      (1786.50 + 998.41)^2 = 2784.91^2 \approx 7{,}755{,}724
    \]

    分母：
    \[
      \frac{1786.50^2}{36} + \frac{998.41^2}{42}
      = \frac{3{,}191{,}582}{36} + \frac{996{,}823}{42}
      \approx 88{,}655 + 23{,}734 = 112{,}389
    \]

    したがって
    \[
      \nu \approx \frac{7{,}755{,}724}{112{,}389} \approx 69.01
    \]

    自由度はおよそ$69.01$.
    端数が出るのが普通である.

    あとは, t統計量 $t \approx 2.28$ を, 自由度およそ$69.01$のt分布と比べるだけだ.
    片側検定の$p$値は, この分布で $t \geq 2.28$ となる確率だから,
    \[
      p = P(T_{69.01} \geq 2.28) \approx 0.013
    \]
    である.

    有意水準を $\alpha = 0.05$ と決めていたとすると, $0.013 < 0.05$ なので帰無仮説は棄却される.
    「原作あり群のほうが平均視聴回数が高い」という主張が, 統計的に支持されたことになる.

    \subsubsection*{なぜ正規分布ではなくt分布なのか}

    ここで使った分布は, 正規分布ではなくt分布だった.
    違いが生じる理由は, 分母に標本から計算した分散が入っている点にある.
    もし母分散がわかっていれば, 分子だけを標準化して正規分布で扱えばよかった.
    ところが母分散は未知で, 手元の標本分散で代用するしかない.
    その代用品自体もばらつく.
    分子と分母が同時にばらつくぶん, 全体としての裾は正規分布より重くなる.
    その重さを正面から扱った分布がt分布だ.

    自由度が大きくなると, 標本分散のばらつきが小さくなり, t分布は正規分布に近づいていく.
    だから標本が大きいときは, t分布で計算しても正規分布で計算しても, $p$値はほとんど変わらない.
    逆に標本が小さいときほど, 正規分布で済ませると, 裾の重さを無視して「実際より珍しい」と判定してしまう.
    小さな標本でこそ, t分布を使う意味がある.


  \subsection{Welchのt検定の自由度はなぜ整数でないのか}

    前節で計算した自由度 $\nu \approx 69.01$ には, 端数がついていた.
    7.2節で見た自由度は, 「集約の代償」であれ「残った自由さ」であれ, データを何個分使ったかという勘定だから, 整数で出てくるはずの量だった.
    個数のはずのものに, なぜ端数がつくのか.

    実は, ここで使われている「自由度」は, 7.2節のものとは少し別の意味になっている.
    どういう経緯で意味が変わったのか, 順を追って見ていく.

    ただ, この話は単に「Welchのt検定の式の都合」では終わらない.
    背景には, 統計的検定そのものに関わる問いがある.
    検定は, 母集団のことが直接わからないからこそ行う作業のはずだ.
    ところが古典的なt検定は, その「わからない」はずの母集団に, いくつかの性質を仮定して組み立てられている.
    その仮定は本当に必要なのか. できるだけ外せないのか.
    端数の自由度は, この問い直しの先に現れる.

    \subsubsection*{ゴセットがしたこと}

    まずは, 古典的なt検定がどう組み立てられたのかを見ておく.
    考案したのは, 6.4節のコラムに登場したゴセットである.
    平均の差をばらつきで割って標準化するというアイデア自体はわかりやすい.
    しかし, その「ばらつき」に手元の標本分散を使うと, 分子と分母が連動してずれる.
    平均の差が偶然大きく出た標本では, 分母のばらつきもまた偶然大きく出やすい, という具合だ.
    この連動を正面から扱うために考え出されたのが, t分布だった.

    ただし, この構成が成り立つように, ゴセットは一つ仮定を置いた.
    \textbf{両群の母分散が等しい}という仮定である.
    この仮定があれば, 二つの群のばらつきを一つにまとめて扱える.
    そして, 自由度 $n_1 + n_2 - 2$ という整数値をもつt分布が出てくる.
    7.2節の「使える情報の個数」という解釈が, ここではきれいに整数のまま生きている.
    一見, 話はうまくまとまっている.

    しかし, そもそも検定をしているのは, 母集団のことが直接わからないからだった.
    その「わからない」はずの母集団について, 「両群で分散は等しい」という性質を仮定してしまっている.
    しかも, 検定の枠組みの中には, この仮定が現実に成り立つかどうかを確かめる手立てがない
    \footnote{
      事前に分散が等しいかを別の検定（F検定など）で調べる, という発想もあるが,
      そうすると「分散が等しいかの検定」と「平均が等しいかの検定」を二段重ねにすることになり,
      別の問題（多重検定）が生じる. かなり厄介な話になっていく.
    }.
    検定の動機と仮定とが, 噛み合っていない感じがする.

    \subsubsection*{仮定を外したい. すると何が起きるか}

    ならば, 等分散の仮定は外したい.
    この方針で組み直したのが, 前節で計算したWelchのt検定である.
    分母に
    \[
      \sqrt{\frac{s_1^2}{n_1} + \frac{s_2^2}{n_2}}
    \]
    を採るのは, 各群のばらつきを最後まで別々に扱うためだ.
    5章のケース3で差に幅をつけたときと同じ形にあたる.

    ところが, 仮定を外した瞬間に厄介なことが起きる.
    分母の中身 $s_1^2/n_1 + s_2^2/n_2$ は, 二つの独立な標本分散を, それぞれ違う重みで足し合わせたものだ.
    両群の分散が等しい場合は, この重みの違いがうまく打ち消し合って, 全体として整数自由度のカイ二乗分布の枠に収まる.
    ゴセットの構成は, このおかげで成立していた.
    分散が違うと, 重み付きの和は整数自由度のカイ二乗分布にはならない.
    そして, 比である $t$ も, 厳密にはt分布に従わなくなる.

    仮定を一つ外しただけのつもりが, t分布で扱うための足場が一気に崩れてしまうのだ.

    \subsubsection*{t分布のままで読めないか}

    ここで, まったく新しい分布を作り直すことも考えられるが, それはあまりにも高コストだ.
    新しい分布の性質を調べ直し, 読み方も一から覚え直すことになる.
    せっかくt分布という慣れた枠組みがあるのだから, 厳密にはt分布でないとしても, 近似的にt分布として読めれば実用上はそれで足りる.

    どうにかしてt分布に寄せたい.
    だが, 何を手がかりに寄せればよいのか.

    手がかりは, 4章の正規分布にある.
    正規分布は, 期待値と分散の二つだけで形が完全に決まる, 際立った分布だった.
    ふつうの分布は, 期待値と分散が同じでも形がいくらでも違いうるのに, 正規分布だけはそうならない.
    この経験を引きずって考えれば, 「分布どうしを寄せたいなら, とりあえず期待値と分散を合わせてみる」というのは, それなりに筋が通って見える.
    この方針が一般の分布で効く保証はないが, 正規分布まわりの素直な分布族であれば, 期待値と分散を揃えるだけでかなり近いところに行けそうだ.
    試す価値はある.

    実際, この見当は当たる.
    重み付きカイ二乗の和を, 一つのカイ二乗分布で「期待値と分散が一致するように」近似すると, 元の分布によく近づくことが知られている.
    この近似がSatterthwaiteのアイデアだ
    \footnote{
      厳密には, 重み付きカイ二乗の和を「定数倍されたカイ二乗分布」で近似している.
      定数倍と自由度の二つを, 期待値と分散の二つの条件で決める, というのが正確な手順だ.
      表に出てくる自由度 $\nu$ は, そのうちの一つにあたる.
    }.
    そして, この近似先のカイ二乗分布の自由度として出てくる値が, 前節で計算した
    \[
      \nu = \frac{\left( \dfrac{s_1^2}{n_1} + \dfrac{s_2^2}{n_2} \right)^2}
                  {\dfrac{(s_1^2/n_1)^2}{n_1 - 1} + \dfrac{(s_2^2/n_2)^2}{n_2 - 1}}
    \]
    だ.

    \subsubsection*{出てきた式を眺めてみる}

    各群のばらつきの寄与 $w_i = s_i^2/n_i$ を主役にして, 自由度の式を書き直してみると,
    \[
      \nu = \frac{(w_1 + w_2)^2}{\dfrac{w_1^2}{n_1 - 1} + \dfrac{w_2^2}{n_2 - 1}}
    \]
    という形になる.
    分母では, 各群の自由度 $n_1 - 1$ と $n_2 - 1$ が, それぞれの群のばらつきの寄与 $w_i$ で重み付けされている.
    全体としては, 二つの群の自由度を, ばらつきの寄与に応じて混ぜ合わせた値が出てきている.
    形としては, 調和平均を一段重み付きにしたものにあたる
    \footnote{
      二つの数を「平均的に」一つにまとめる操作に, 算術平均ではなく調和平均的な形が出てくるのは少し意外かもしれない.
      $\nu$ の式を二乗和の構造から眺めると自然に出てくるものだが, 本書ではこれ以上は立ち入らない.
    }.

    両群のばらつきの寄与が等しい, つまり $w_1 = w_2$ の場合, $\nu$ は $n_1 - 1$ と $n_2 - 1$ の調和平均の二倍にあたる値になる.
    さらに $n_1 = n_2$ なら, ちょうど $n_1 + n_2 - 2$ となり, 古典的なt検定の自由度と一致する.

    一方, ばらつきの寄与に偏りがあったり件数が違ったりすると, $\nu$ は整数にならない.
    重みも件数の比も実数値だから, それらを混ぜ合わせた結果が整数になる必然性はない.
    端数が出るのが普通なのだ.

    \subsubsection*{t分布族はもともと実数で定義されている}

    ここで一つ, 引っかかる点が残る.
    端数の自由度をt分布に「入れる」とは, どういう操作なのか.
    整数のときの定義を, 実数の場合に勝手に拡張していいのか.

    実は, 拡張という言い方は正確ではない.
    t分布族は, 最初から実数のパラメータで定義された分布族だからだ.
    密度関数を書くと
    \[
      f(t; \nu) = \frac{\Gamma\!\left(\dfrac{\nu+1}{2}\right)}
                       {\sqrt{\nu\pi}\, \Gamma\!\left(\dfrac{\nu}{2}\right)}
                  \left(1 + \frac{t^2}{\nu}\right)^{-\frac{\nu+1}{2}}
    \]
    のようになる.
    式の細部は本書の範囲を超えるが, 見たいのはパラメータ $\nu$ の使われ方だけである.
    $\nu$ は分母の根号の中や指数の肩にも入っており, さらに, 階乗を実数へ滑らかに延ばした関数（ガンマ関数 $\Gamma$）を介して使われている.
    その結果, $\nu$ には正の実数なら何を入れても, 密度関数として意味を持つ.
    積分すれば必ず1になるし, 確率分布として一貫している.

    つまりt分布族は, 整数のために作られた分布族を後から実数に拡張したものではない.
    最初から, 実数のパラメータで連続的に張られた分布族として存在している.
    その中で, $\nu$ が整数のときに限って「独立な情報の個数」という統計的解釈が乗る.
    整数自由度のt分布は, 連続な分布族の中の, とびとびの特別な点にすぎなかったのだ.

    だから $\nu \approx 69.01$ のような実数値を入れることは, t分布族にもともと備わっている性質を使っているだけのことで, 何ら不自然なことは起きていない.

    \subsubsection*{「自由度」という言葉が二つの意味を持っている}

    ここまでを並べると, 「自由度」という言葉が二つの意味で使われていることがわかる.
    7.2節の自由度は, 「集約の代償」であり「残った自由さの個数」だった.
    どちらの見方でも, 出てくる数は整数である.
    一方, Welchのt検定の自由度は, t分布族の中で近似先の位置を指定するパラメータの値である.
    こちらは代償でも個数でもない.
    整数の場合に限って7.2節の解釈が乗るので, その名前を実数の場合にまで流用しているにすぎない.

    この区別を知らずに「自由度が69.01だった」と聞くと, 「集約の代償が69.01個分とはどういうことか」と混乱する.
    実際に起きているのは, 概念の拡張である.
    7.2節で掴んだ「集約の代償」という解釈は, 整数自由度の世界では今も有効だ.
    Welchの自由度では, 同じ言葉が「t分布族の中の位置」という別の意味で使われている.

    \subsubsection*{Welchで何を仮定しているか, そして仮定が崩れたら}

    話を仮定の整理に戻したい.
    Welchのt検定が「両群の母分散が等しい」という仮定を外した, ということを見てきた.
    では, 残されている仮定は何か.

    まず必要なのは, 観測の独立性である.
    どの観測も, 他の観測に影響されずに得られている, という前提だ.
    この仮定を外す話は本書の範囲を超えるので, ここでは置く.

    次に, 各群の中では, 観測値がそれぞれ同じ母集団分布から得られていると考える.
    つまり, 原作あり群の中の観測値は一つの分布から, 原作なし群の中の観測値もまた別の一つの分布から得られている, という見方をする.
    Welchのt検定が外しているのは, この二つの分布が同じでなければならない, という仮定のほうだ.
    群内では同じ, 群間では違ってよい, という整理になる.

    そして, t統計量の分子である $\bar{x}_1 - \bar{x}_2$ が正規分布で扱えること.
    古典的な導出では, 各群の母集団がそれぞれ正規分布に従うとし, 独立な正規分布に従う変数の差はやはり正規分布になることを使う.
    この限りでは「両群がそれぞれ正規分布に従う」という強めの条件が要る.
    ただし実用上は, もっと緩くて済む.
    標本サイズがそれなりに大きければ, 4章の中心極限定理によって, 各群の標本平均は近似的に正規分布に従う.
    独立な近似正規変数の差もまた近似的に正規になるから, 結局 $\bar{x}_1 - \bar{x}_2$ も近似的に正規分布で扱える.
    つまり, $n$ が大きければ, 元の母集団が完全に正規でなくても, この仮定はほぼ気にしなくてよくなる.

    Studentのt検定とWelchのt検定の関係を, ここで一度整理しておこう.
    Studentでは, 「両群は同じ広がりをもち, 平均だけが異なるかもしれない正規分布」を考える.
    Welchでは, 「平均も広がりも群ごとに異なってよい」と考える.
    どちらも, 「群内では同じ分布から独立に得られている」という前提は共通で, Welchが緩めたのは群間の広がりに関する制約だ.

    したがって実用上は, 標本サイズがそれなりにあって, 観測が独立で, 群内のデータの出どころが大きく混ざっていなければ, 元の母集団がきれいな正規分布でなくてもWelchのt検定は使いやすい.
    一方で, 標本サイズが小さく, かつ元の母集団が大きく歪んでいる場合や, 外れ値が平均を強く動かしている場合には, この近似が効かなくなり, Welchの結果も信頼しにくくなる.
    そういう場合には別の手立てを探すことになる.
    たとえば, 平均そのものを比べるのではなく順位だけを比べる方法（順位検定）や, 5章でも名前だけ触れた, 手元のデータからリサンプリングを繰り返して幅をつかむ方法（ブートストラップ）がある.
    本書ではこれ以上は扱わないが, 必要になったときに調べられるよう, 名前を残しておく.


  \subsection{カイ二乗検定の数理構造と計算の実際}

    Part Aでは, カテゴリデータの検定として, 「血液型と性格タイプに関係があるのか」という問いを取り上げた.
    ここではその例に戻って, カイ二乗検定の計算を実際に進めながら, 統計量と自由度の意味を確かめる.

    100人に血液型と性格タイプを尋ねた結果は, 次の通りだった：

    \begin{center}
      \begin{tabular}{l|cccc|c}
            & A型 & B型 & O型 & AB型 & 合計 \\
      \hline
      几帳面 & 18  & 6   & 10  & 8    & 42 \\
      大雑把 & 7   & 19  & 15  & 17   & 58 \\
      \hline
      合計   & 25  & 25  & 25  & 25   & 100
      \end{tabular}
    \end{center}

    仮説は次のように立てる.
    \begin{itemize}
      \item 帰無仮説 $H_0$：血液型と性格タイプは無関係（独立）である
      \item 対立仮説 $H_1$：血液型と性格タイプには何らかの関係がある（独立ではない）
    \end{itemize}
    特定の方向を仮定していないので, 両側検定となる.

    \subsubsection*{期待度数の計算}

    帰無仮説のもとでは, 血液型と性格タイプは独立である.
    このとき各マスに入るはずの人数, つまり期待度数は, 次のように計算できるのだった.

    \[
    \text{期待度数} = \frac{\text{行の合計} \times \text{列の合計}}{\text{全体の合計}}
    \]

    たとえば, 「A型で几帳面」の期待度数は
    \[
    \frac{42 \times 25}{100} = 10.5
    \]
    となる.
    同様に他のマスの期待度数を計算すると, 次のようになる：

    \begin{center}
      \begin{tabular}{l|cccc}
            & A型 & B型 & O型 & AB型 \\
      \hline
      几帳面（期待） & 10.5 & 10.5 & 10.5 & 10.5 \\
      大雑把（期待） & 14.5 & 14.5 & 14.5 & 14.5
      \end{tabular}
    \end{center}

    \subsubsection*{カイ二乗統計量の計算}

    カイ二乗統計量 $\chi^2$ は, 観測度数と期待度数のずれの大きさを, 次の式で集計する：

    \[
    \chi^2 = \sum_{\text{すべてのマス}} \frac{(O - E)^2}{E}
    \]

    ここで $O$ は観測度数, $E$ は期待度数である.
    ずれの二乗 $(O - E)^2$ をそのまま足さずに $E$ で割っているのは, ずれの重みを揃えるためだ.
    期待度数が100人のマスで5人ずれるのと, 期待度数が10人のマスで5人ずれるのとでは, 珍しさがまるで違う.
    各マスのずれを, そのマスで見込まれる規模に対する相対的な大きさに直してから, 足し合わせている.

    各マスについてこの値を計算してみよう（小数第3位で四捨五入）：

    \begin{itemize}
      \item $(18 - 10.5)^2 / 10.5 = 56.25 / 10.5 \approx 5.36$
      \item $(6 - 10.5)^2 / 10.5 = 20.25 / 10.5 \approx 1.93$
      \item $(10 - 10.5)^2 / 10.5 = 0.25 / 10.5 \approx 0.02$
      \item $(8 - 10.5)^2 / 10.5 = 6.25 / 10.5 \approx 0.60$
      \item $(7 - 14.5)^2 / 14.5 = 56.25 / 14.5 \approx 3.88$
      \item $(19 - 14.5)^2 / 14.5 = 20.25 / 14.5 \approx 1.40$
      \item $(15 - 14.5)^2 / 14.5 = 0.25 / 14.5 \approx 0.02$
      \item $(17 - 14.5)^2 / 14.5 = 6.25 / 14.5 \approx 0.43$
    \end{itemize}

    これらをすべて足し合わせると,

    \[
    \chi^2 \approx 5.36 + 1.93 + 0.02 + 0.60 + 3.88 + 1.40 + 0.02 + 0.43 = 13.64
    \]

    \subsubsection*{自由度と\texorpdfstring{$p$}{p}値}

    この統計量を比べる先は, カイ二乗分布である.
    そしてカイ二乗分布の形も, t分布と同じく自由度で決まる.
    クロス集計表の検定では, 自由度は
    \[
    (\text{行数} - 1) \times (\text{列数} - 1) = (2 - 1) \times (4 - 1) = 3
    \]
    である.

    なぜ行数や列数から1を引くのか.
    期待度数を作るとき, 行の合計と列の合計を標本から計算して使った.
    合計が固定されると, 表のマスを自由に埋められる余地が減る.
    実際, この2行4列の表では, たとえば「几帳面」の行のA型・B型・O型の3マスを決めると, 残りのマスはすべて行と列の合計から自動的に決まってしまう.
    自由に決められるマスの数は $(2-1) \times (4-1) = 3$ 個しかない.
    7.2節で見た「平均を固定すると, 最後の一つは帳尻合わせで決まる」という勘定と同じことが, ここでは行と列の両方向で起きている.

    したがって, $\chi^2 = 13.64$ を比べる先は, 自由度3のカイ二乗分布になる.
    この分布で13.64以上の値が出る確率を求めると,
    \[
    p = P(\chi^2_3 \geq 13.64) \approx 0.0035
    \]
    である.

    有意水準を $0.05$ としていたとすると, $p$値はそれよりはるかに小さい.
    帰無仮説は棄却され, 「血液型と性格タイプは無関係とは言いがたい」という結論になる.
    もっとも, この結論をどこまで信じてよいかについては, Part Aで触れたとおり, データの集め方の問題が残っている.

    \subsubsection*{なぜカイ二乗分布と比べられるのか}

    カイ二乗分布は, 理論的には, 標準正規分布に従う独立な変数を二乗して足し合わせたものが従う分布として定義され, 足し合わせる個数が自由度になる.
    観測度数と期待度数のずれは, 標本が大きければ近似的に正規分布に従うことが知られており, その二乗和であるカイ二乗統計量は, 自由に動ける個数分, つまり自由度3のカイ二乗分布で近似できる.
    この導出の細部は本書の範囲を超えるが, 「統計量を, 帰無仮説のもとでの分布と比べて珍しさを測る」という枠組み自体は, t検定とまったく同じである.


  \subsection{検定の二つの誤りと限界}

    ここまで, t検定とカイ二乗検定の計算を通じて, データの差やずれが偶然で説明できるかどうかを判断してきた.
    しかし, どれだけ丁寧に手続きを踏んでも, 検定の結論が間違っている可能性は消せない.
    Part Aの新薬のコラムで名前だけ触れた「二つの誤り」を, ここで定式化しておこう.

    \subsubsection*{検定の結果は4通りに分かれる}

    帰無仮説が実際に正しいかどうかと, 検定がそれを棄却するかどうかの組み合わせで, 検定の結果は次の4通りに分かれる：

    \begin{center}
      \begin{tabular}{c|cc}
            & \textbf{帰無仮説が正しい} & \textbf{帰無仮説が誤っている} \\
      \hline
      \textbf{棄却しない} & 正しい判断 & \textbf{第II種の誤り（見逃し）} \\
      \textbf{棄却する}   & \textbf{第I種の誤り（誤検出）} & 正しい判断 \\
      \end{tabular}
    \end{center}

    統計的検定では, このうち第I種の誤りを抑えることをまず優先する.
    その上限を決めているのが, 有意水準 $\alpha$ である.

    \subsubsection*{第I種の誤り（false positive）}

    第I種の誤りとは, 実際には差がないのに「統計的に有意」と判断してしまう誤りである.
    有意水準を $\alpha = 0.05$ に設定するというのは, 差がない世界で検定を行ったときに, 偶然だけで有意という結果が出る確率を5\%以下に抑える, という設計を意味する.

    裏を返すと, $p$値が有意水準を下回ったとしても, 「差がない世界で, たまたま極端な結果を引き当てた」可能性そのものは消えていない.
    有意という結果から「必ず差がある」とまでは言えないのは, このためである.

    \subsubsection*{第II種の誤り（false negative）}

    第II種の誤りとは, 実際には差があるのに検定でそれを検出できず, 「差があるとは言えない」と判断してしまう誤りである.

    この誤りは, サンプルサイズが小さい場合に起こりやすい.
    ばらつきの大きいデータでは, 本当に平均に差があっても, その差が偶然の変動に埋もれて, t統計量があまり大きくならないことがあるからだ.

    第II種の誤りが起こる確率は $\beta$ と書かれることが多い.
    その補数 $1 - \beta$, すなわち「本当に差があるときに, それを検出できる確率」を\textbf{検出力}（statistical power）と呼ぶ.

    \subsubsection*{\texorpdfstring{$p$}{p}値を決めているのは差の大きさだけではない}

    二つの誤りの背景には, サンプルサイズの影響がある.
    サンプルサイズが大きいほど, 平均のばらつき（標準誤差）は小さくなり, t統計量は大きく出やすい.
    その結果, ごく小さな差であっても有意になることがある.
    逆にサンプルサイズが小さいと, 統計量が有意水準に届かず, 差を見逃す（第II種の誤りを犯す）可能性が高くなる.

    つまり, $p$値の大小は「差の大きさ」だけで決まるのではなく, 差の大きさ, サンプルサイズ, ばらつきの三つで決まる.
    「$p$値が有意だから差が大きい」とも, 「$p$値が有意でないから差がない」とも, 短絡的に結論することはできない.

    \subsubsection*{検定だけで決めない}

    統計的検定が定量化しているのは, 「観察されたデータが, 帰無仮説とどれくらい食い違っているか」である.
    $p$値が一つ出てきたからといって, それだけで主張の正しさが決まるわけではない.
    検定の結論は, 仮説の立て方（片側か両側か）, 有意水準の設定, サンプルの取り方（偏りや独立性）といった前提の上に載っていて, 前提のどれかが崩れていれば, $p$値の意味も一緒に崩れる.
    検定の結果は, こうした前提とあわせて, 判断材料の一つとして使う.


  \subsection{効果量と実質的な差の検討}

    検定の結論は, 「差が統計的に有意かどうか」という形で出てくる.
    しかし, 有意だとわかったとして, その差が実際に重要な意味を持つかどうかは, 別の問題である.

    たとえば, 大規模な調査で, ある栄養ドリンクを飲んだグループの平均寿命が, 飲んでいないグループよりも0.3日長く, この差が$p$値0.01で有意だったとしよう.
    統計的な手続きとしては, 文句のつけようがない.
    それでも, 寿命が0.3日延びることに, このドリンクを毎日飲み続けるだけの価値があるだろうか.

    このように, 「統計的に有意な差」と「実質的に意味のある差」は, しばしば一致しない.
    そこで, 差の大きさそのものを測る指標である\textbf{効果量}（effect size）を, 検定とあわせて確認する.
    $p$値が「差があると言えるか」に答えるのに対し, 効果量は「その差はどのくらい大きいか」に答える.
    効果量は, 検定で扱った差の大きさを, 単位やサンプルサイズに依存しない形で表した指標である.

    \subsubsection*{t検定における効果量：Cohenの\texorpdfstring{$d$}{d}}

    t検定でもっともよく使われる効果量に, \textbf{Cohenの$d$}がある.
    考え方は素朴で, 二つのグループの平均の差を, 各群のばらつきを代表する標準偏差で割る.

    \[
      d = \frac{\bar{x}_1 - \bar{x}_2}{s}
    \]

    割る前の差は「点」や「回」といった単位を引きずっているが, 同じ単位の標準偏差で割ると, 「ばらつき何個分の差か」という単位のない数に直る.
    だから, テストの点数の差でも視聴回数の差でも, 同じ基準で大きさを比べられる.
    ここで $s$ は両群のばらつきを代表する標準偏差で, 各群の分散の平均をとってその平方根を採るのが本書の立場である：

    \[
      s = \sqrt{\frac{s_1^2 + s_2^2}{2}}
    \]

    原作あり／なしの例で計算してみよう.
    \[
      s = \sqrt{\frac{257.1^2 + 207.2^2}{2}} = \sqrt{\frac{66100.41 + 42931.84}{2}} = \sqrt{54516.13} \approx 233.49
    \]

    したがって
    \[
      d = \frac{274.1 - 153.7}{233.49} \approx 0.52
    \]

    Cohenの$d$は, およそ$0.52$となる.

    \subsubsection*{効果量の目安}

    Cohenは, $d$の大きさの目安として以下を提案している：

    \begin{itemize}
      \item $d \approx 0.2$ ： 小さい効果（small）
      \item $d \approx 0.5$ ： 中程度の効果（medium）
      \item $d \approx 0.8$ ： 大きな効果（large）
    \end{itemize}

    上の例の $d \approx 0.52$ は, この目安では「中程度の効果」にあたる.
    平均の差120.4万回は, 両群のばらつきを代表する標準偏差（約233万回）のおよそ半分にあたる, という読みになる.
    もちろん, この区切りも有意水準の0.05と同じで, 絶対的な根拠があるわけではない.
    それでも, 分野を越えて差の大きさを伝え合うための, 共通の目安になっている.

    \subsubsection*{\texorpdfstring{$p$}{p}値と効果量はセットで見る}

    $p$値は「差があると言えるか」を受け持ち, 効果量は「その差がどのくらい大きいか」を受け持つ.
    どちらか一方だけでは, 検定の結果を読み違える.
    サンプルサイズが巨大なら, 効果量の小さな差でも$p$値は有意になるし, サンプルサイズが小さければ, 効果量の大きな差でも有意にならないことがある.
    だから近年の研究報告では, 有意かどうかだけでなく, 効果量を併記することが強く求められている.

    Part Bでは, Part Aで概念として追った検定の枠組みを, 計算で確かめてきた.
    検定統計量と自由度が決まれば$p$値が決まり, $p$値と有意水準の比較で判断が下りる.
    ただし, その判断には第I種・第II種の誤りの可能性が常に残り, 差の実質的な大きさには効果量が別に答える.
    $p$値と効果量とサンプルサイズをセットで読むところまでが, 検定を使うということである.
